\chapter*{Glossary}
\addcontentsline{toc}{chapter}{Glossary}
\label{chap:glossary}

This glossary explains the recurring technical terms of the thesis.\kiref{Claude Opus 5/3}
Terms from machine-learning practice are defined in the sense in which they are
standard in that literature; where this thesis defines a term of its own, the entry
says so. The formal equations for the evaluation metrics are given in
Appendix~\ref{sec:method-metric-definitions}.

{\renewcommand{\arraystretch}{1.35}
\begin{longtable}{p{0.27\textwidth} p{0.66\textwidth}}
\hline
\textbf{Term} & \textbf{Definition} \\
\hline
\endfirsthead
\multicolumn{2}{l}{\small\itshape Glossary --- continued from previous page} \\
\hline
\textbf{Term} & \textbf{Definition} \\
\hline
\endhead

Ablation &
An experiment that removes or disables one component of a system and retrains it,
so that the difference in results measures what that component actually
contributed. Used here for the human-actor filter
(Appendix~\ref{appendix:bot-ablation}) and for feature availability
(Section~\ref{sec:method-prediction-regimes}). \\

Average precision &
Area under the precision-recall curve: a ranking measure that, unlike ROC-AUC,
ignores correctly ranked negatives and so is not flattered by a large negative
class. Its no-skill reference is the positive rate of the evaluated set, not
\(0.5\), so it must be read alongside that rate
(Section~\ref{sec:results-precision-recall}). \\

Brier score &
A measure of how accurate predicted \emph{probabilities} are: the mean squared
difference between the predicted probability and the observed outcome. Lower is
better. \\

Calibration &
Whether predicted probabilities match reality in the long run --- if the model
says \(0.8\) for a group of repositories, roughly \(80\%\) of them should in fact
become inactive. Distinct from discrimination: a model can rank correctly and
still be badly calibrated. \\

Cold-start &
The standard machine-learning term for having to make a prediction about an entity
for which no history is available. Here it names the model that uses only the 17
current-snapshot features, and applies to any repository absent from the
precomputed feature cache. \\

Discrimination &
Whether the model puts the repositories that became inactive above those that did
not, irrespective of the absolute probabilities. Measured by ROC-AUC. \\

Expected calibration error (ECE) &
A single number summarizing miscalibration: predictions are grouped into
probability bins and the average gap between each bin's predicted probability and
its observed outcome frequency is reported. \\

F1 score &
The harmonic mean of precision and recall at a chosen decision threshold, i.e.\ a
single number balancing how many flagged repositories were really inactive against
how many inactive repositories were flagged. Unlike ROC-AUC it depends on the
threshold. \\

Full-history &
The counterpart to cold-start: the model that additionally uses the 43 features
reconstructed from the event history before the observation date. \\

Held-out test set &
The rows set aside before training and used neither for fitting the model nor for
calibrating its probabilities, so that the reported metrics describe performance
on data the model has never seen. \\

Human-actor filter &
A filter of this thesis that excludes bot accounts from commit- and
contributor-derived signals and from the label, so that automated activity cannot
be mistaken for maintenance (Section~\ref{sec:method-feature-engineering}). \\

Inactivity label &
The prediction target: \(1\) if the repository is \emph{not} maintained during the
twelve months after the observation date, \(0\) otherwise
(Section~\ref{sec:method-definitions}). \\

Leakage &
The error of letting information that would not have been available at prediction
time influence the model, which inflates measured performance. Prevented here by
the temporal boundary at the observation date
(Section~\ref{sec:method-leakage-prevention}). \\

Observation snapshot &
One row of the dataset: a single repository described as it appeared at one
observation date, together with the label derived from the following twelve
months. \\

Prediction regime &
The term used in this thesis for which of the two nested feature sets a model was
trained on --- full-history or cold-start. \\

Reconstructed proxy &
A feature that measures the same underlying construct as an established metric but
is computed by this thesis from GH Archive event data under its own rules, and is
therefore not that metric. The responsiveness features are reconstructed proxies
for, not implementations of, the corresponding CHAOSS metrics
(Section~\ref{sec:related-responsiveness}). \\

Right-censored &
A value that is known only as a lower bound. If a repository has no observed
commit, ``days since last commit'' has no exact value but is at least the length
of the observable record, and is encoded at that bound rather than as zero
(Section~\ref{sec:method-feature-engineering}). \\

ROC-AUC &
Area under the receiver operating characteristic curve: the probability that a
randomly chosen inactive repository is scored above a randomly chosen active one.
\(0.5\) is a coin flip, \(1.0\) a perfect ranking. \\

Slice evaluation &
Re-computing the metrics separately within meaningful subgroups of the test set
(for example by popularity, or restricted to repositories still active at the
observation date), so that strong performance on a large subgroup cannot hide weak
performance on a smaller one. The word \emph{slice} is used in this thesis only for
such a subgroup; the three data partitions are called the training, validation, and
test sets. \\

Two-of-four maintenance rule &
This thesis's own operational definition of maintenance: a repository counts as
maintained in the twelve-month window if at least two of four activity checks are
satisfied (Section~\ref{sec:method-definitions}). It is a design choice of this
work, not an established external standard, and its sensitivity is tested in
Section~\ref{sec:results-label-sensitivity}. \\
\hline
\end{longtable}
}
